% Minimal preamble for Open Logic Project sources (pinned commit 1e960beff9ed7835bf3e3f1335e21af3439cd107).
%
% OLP's own master preamble cannot run on the project host: TeX Live 2025 there lacks
% prooftree.sty, ntheorem.sty, extarrows.sty, mathdots.sty, isomath.sty and subfiles.sty.
% Following the project's own rule ("if the author preamble needs packages the toolchain lacks,
% do not invent macros -- use a minimal working preamble"), this file declares, verbatim in
% spirit and numbering behaviour, exactly the theorem-like environments OLP defines in
% sty/open-logic-book-envs.sty (thm/lem/prop/cor/defn/ex/prob/rem/note/case/axiom/conv),
% plus the shorthand accents OLP uses (`!!a`, `!!b`, `!!{...}`).
%
% Author text is untouched; only the packaging is ours.
\documentclass[10pt,a4paper]{article}

\usepackage{amsmath}
\usepackage{amssymb}
\usepackage{amsthm}
\usepackage{graphicx}
\usepackage{hyperref}
\usepackage{enumitem}
\usepackage[framemethod=default]{mdframed}

% ---- OLP shorthand accents (from sty/open-logic.sty) ----
% NB: \! is taken by LaTeX (negative thin space), so we only define the !!a / !!b families
% and the \olup helper used by them.
\newcommand*{\olup}[1]{\MakeUppercase{#1}}
\def\!!!a#1{\olup{#1}}
\def\!!a#1{\olup{#1}}
\def\!!b#1{\olup{#1}}

% ---- theorem-like environments, same family numbering as OLP ----
\theoremstyle{plain}
\newtheorem{thm}{Theorem}[section]
\newtheorem{lem}[thm]{Lemma}
\newtheorem{prop}[thm]{Proposition}
\newtheorem{cor}[thm]{Corollary}
\newtheorem{axiom}{Axiom}

\theoremstyle{definition}
\newtheorem{defn}[thm]{Definition}
\newtheorem{ex}[thm]{Example}
\newtheorem{prob}{Problem}[section]
\newtheorem{probd}{Problem}
\newtheorem{conv}[thm]{Convention}

\theoremstyle{remark}
\newtheorem{rem}[thm]{Remark}
\newtheorem{note}[thm]{Note}
\newtheorem{case}[thm]{Case}

% ---- environments OLP defines as wrappers ----
\newenvironment{probtag}[1]{\begin{prob}}{\end{prob}}
\newenvironment{explain}{}{}
\newenvironment{explain*}{}{}
\newenvironment{editorial}{}{}
\newenvironment{digress}{}{}
\newenvironment{digress*}{}{}
\newenvironment{history}{}{}
\newenvironment{reference}{}{}
\newcommand*{\olfileid}[3]{}
\newcommand*{\olsection}[1]{\section*{#1}}
\newcommand*{\olchapter}[3]{}
\newcommand*{\olpart}[2]{}
\newcommand*{\olimport}[1]{}
\newcommand*{\olinclude}[1]{}
\newcommand*{\olpath}[1]{}
\newcommand*{\ollangid}[1]{}
\newcommand*{\problemtag}[1]{}
\newcommand*{\problemauthor}[1]{}
\newcommand*{\ollabel}[1]{\label{#1}}
% \olref takes 0-3 OPTIONAL arguments plus one mandatory label:
%     \olref{lab}  |  \olref[from]{lab}  |  \olref[from][to]{lab}  |  \olref[a][b][c]{lab}
% A single-argument definition silently swallows "[sth]" and emits a broken reference
% (measured on the host: the text came out as "??iso]x"). The chain below consumes the
% brackets and cross-references only the mandatory argument.
\makeatletter
\newcommand{\olref}{\@ifnextchar[{\olref@i}{\olref@final}}
\def\olref@i[#1]{\@ifnextchar[{\olref@ii}{\olref@final}}
\def\olref@ii[#1]{\@ifnextchar[{\olref@iii}{\olref@final}}
\def\olref@iii[#1]{\olref@final}
\def\olref@final#1{\ref{#1}}
\makeatother
\newcommand*{\olnewpage}{\clearpage}

% ===================== OLP compatibility layer (1) =====================

\makeatletter

% ---------------- proof trees (OLP's prooftree is a bussproofs derivative) ----------------
\providecommand{\AxiomC}[1]{}
\providecommand{\UnaryInfC}[1]{}
\providecommand{\BinaryInfC}[1]{}
\providecommand{\TrinaryInfC}[1]{}
\providecommand{\QuaternaryInfC}[1]{}
\providecommand{\QuinaryInfC}[1]{}
\providecommand{\RightLabel}[1]{}
\providecommand{\LeftLabel}[1]{}
\providecommand{\doubleLine}{}
\providecommand{\noLine}{}
\providecommand{\DisplayProof}{}
\providecommand{\fCenter}{}
\newenvironment{prooftree}{}{}
\newenvironment{ltprooftree}{}{}
\newenvironment{oltableau}{}{}

% ---------------- logical notation ----------------
\providecommand{\lif}{\mathbin{\rightarrow}}
\providecommand{\lforall}{\mathord{\forall}}
\providecommand{\lexists}{\mathord{\exists}}
\providecommand{\ltrue}{\top}
\providecommand{\lfalse}{\bot}
\providecommand{\Ltrue}{\top}
\providecommand{\Lfalse}{\bot}
\providecommand{\Sequent}{\mathrel{\Rightarrow}}
\providecommand{\Proves}{\mathrel{\vdash}}
\providecommand{\Entails}{\mathrel{\models}}
\providecommand{\Subst}[3]{#1[#2/#3]}
\providecommand{\sFmla}[2]{#1(#2)}
\providecommand{\pFmla}[2]{#1(#2)}
\providecommand{\formula}{\relax}
\providecommand{\True}{\top}
\providecommand{\False}{\bot}
\providecommand{\Args}[1]{#1}
\providecommand{\Ts}{\ensuremath{\top}}
\providecommand{\Fi}{\ensuremath{\bot}}
\providecommand{\lfalsum}{\bot}

% ---------------- objects, structures, metatheory ----------------
\providecommand{\Obj}{}
\providecommand{\Struct}{\ensuremath{\mathfrak{A}}}
\providecommand{\Domain}{\ensuremath{|\mathfrak{A}|}}
\providecommand{\Assign}{\ensuremath{s}}
\providecommand{\Sat}[2]{\ensuremath{#1 \models #2}}
\providecommand{\mSat}[3]{\ensuremath{#1 \models #2[#3]}}
\providecommand{\mModel}[2]{\ensuremath{#1 \models #2}}
\providecommand{\Lang}{\ensuremath{\mathcal{L}}}
\providecommand{\Th}{\ensuremath{\mathrm{Th}}}
\providecommand{\Log}{\ensuremath{\mathrm{L}}}
\providecommand{\tuple}[1]{\langle #1 \rangle}
\providecommand{\num}[1]{\underline{#1}}
\providecommand{\cfind}[1]{\varphi_{#1}}
\providecommand{\Nat}{\ensuremath{\mathbb{N}}}
\providecommand{\Mult}{\ensuremath{\mathrm{Mult}}}
\providecommand{\Add}{\ensuremath{\mathrm{Add}}}
\providecommand{\Exp}{\ensuremath{\mathrm{Exp}}}
\providecommand{\Zero}{\ensuremath{0}}
\providecommand{\Succ}{\ensuremath{\mathrm{S}}}
\providecommand{\lambd}{\ensuremath{\lambda}}
\providecommand{\bredpar}{\mathrel{\rightarrow_\beta}}
\providecommand{\bcd}[1]{\ensuremath{\ulcorner #1 \urcorner}}
\providecommand{\Set}[1]{\ensuremath{\{#1\}}}
\providecommand{\indcase}[1]{\par\noindent\emph{Case #1}.\quad}
\providecommand{\indfrm}[1]{\par\noindent\emph{Form #1}.\quad}
\providecommand{\iftag}[1]{#1}
\providecommand{\tagitem}[1]{\item[#1]}
\providecommand{\Elim}[1]{}
\providecommand{\Intro}[1]{}
\providecommand{\LeftR}[1]{}
\providecommand{\RightR}[1]{}
\providecommand{\TAss}[1]{}
\providecommand{\TRule}[1]{}

\makeatother
\providecommand{\olpath}{.}

% ===================== OLP compatibility layer (2) =====================
% --- second compatibility block: notation and shorthands used by the frozen OLP units ---
\providecommand{\CutCS}{\mathrm{Cut}}
\providecommand{\LeftSubproofLabel}{\relax}
\providecommand{\RightSubproofLabel}{\relax}
\providecommand{\maeh}{\relax}
\providecommand{\cheight}{\relax}
\providecommand{\ande}{\mathrm{and}_E}
\providecommand{\andi}{\mathrm{and}_I}
\providecommand{\ore}{\mathrm{or}_E}
\providecommand{\ori}{\mathrm{or}_I}
\providecommand{\typeof}{\mathord{:}}
\providecommand{\pValue}[1]{#1}
\providecommand{\Sent}{\mathrm{Sent}}
\providecommand{\cardplus}{\mathbin{+}}
\providecommand{\cardtimes}{\mathbin{\cdot}}
\providecommand{\cardexpo}{\mathbin{\hat{\ }}}
\providecommand{\cardnless}{\not<}
\providecommand{\cardle}{\leq}
\providecommand{\funfromto}[3]{#1\colon #2 \to #3}
\providecommand{\funimage}[1]{\operatorname{im}(#1)}
\providecommand{\fdefined}{\mathord{\downarrow}}
\providecommand{\canonord}{\operatorname{ord}}
\providecommand{\ordtimes}{\mathbin{\cdot}}
\providecommand{\ordeq}[2]{#1 = #2}
\providecommand{\ordneq}[2]{#1 \neq #2}
\providecommand{\ordtype}[1]{\operatorname{ord}(#1)}
\providecommand{\ordsucc}{\mathord{+1}}
\providecommand{\ordunion}{\cup}
\providecommand{\onesphere}{S^1}
\providecommand{\rotationsgroup}{\mathrm{SO}(2)}
\providecommand{\equivrep}[1]{[#1]}
\providecommand{\rlexless}{\prec}
\providecommand{\ZFC}{\mathrm{ZFC}}
\providecommand{\ZF}{\mathrm{ZF}}
\providecommand{\Rat}{\mathbb{Q}}
\providecommand{\Real}{\mathbb{R}}
\providecommand{\PosInt}{\mathbb{Z}^{+}}
\providecommand{\Id}{\mathrm{id}}
\providecommand{\Bin}{\mathrm{Bin}}
\providecommand{\pto}{\rightharpoonup}
\providecommand{\dom}{\operatorname{dom}}
\providecommand{\ran}{\operatorname{ran}}
\providecommand{\len}{\operatorname{len}}
\providecommand{\concat}{\mathbin{{}^\frown}}
\providecommand{\card}[1]{\lvert #1 \rvert}
\providecommand{\Pow}{\mathcal{P}}
\providecommand{\Setabs}[2]{\{\, #1 \mid #2 \,\}}
\providecommand{\comp}{\circ}
\providecommand{\Char}{\chi}
\providecommand{\Prf}{\operatorname{Prf}}
\providecommand{\Prov}{\operatorname{Prov}}
\providecommand{\Proj}{\operatorname{proj}}
\providecommand{\Gn}{\operatorname{gn}}
\providecommand{\gn}{\operatorname{gn}}
\providecommand{\Value}[1]{\operatorname{val}(#1)}
\providecommand{\Deduce}{\mathrm{D}}
\providecommand{\Dual}{\mathrm{d}}
\providecommand{\Ax}{\mathrm{Ax}}
\providecommand{\pSat}[3]{#1 \models #2\,[#3]}
\providecommand{\pAssign}{\sigma}
\providecommand{\bforall}{\forall}
\providecommand{\bexists}{\exists}
\providecommand{\liff}{\mathbin{\leftrightarrow}}
\providecommand{\FV}{\operatorname{FV}}
\providecommand{\UnaryInf}{\ensuremath{\mid}}
\providecommand{\BinaryInf}{\ensuremath{\mid}}
\providecommand{\depth}{\operatorname{depth}}
\providecommand{\oliflabeldef}[3]{#1}
\providecommand{\TMblank}{\sqcup}
\providecommand{\TMstroke}{\mid}
\providecommand{\TMstay}{\mathrm{S}}
\providecommand{\TMright}{\mathrm{R}}
\providecommand{\TMleft}{\mathrm{L}}
\providecommand{\TMendtape}{\triangleright}
\providecommand{\umin}{\mu}
\providecommand{\Sum}{\sum}
\providecommand{\citep}[1]{\cite{#1}}
\providecommand{\citet}[1]{\cite{#1}}
\providecommand{\citeauthor}[1]{\cite{#1}}
\providecommand{\citeyear}[1]{\cite{#1}}

% ===================== OLP compatibility layer (3) =====================
% --- third compatibility block: remaining notation seen in the frozen OLP units ---
\providecommand{\iso}{\cong}
\providecommand{\nszero}{\mathrm{0}}
\providecommand{\nssucc}{\mathrm{S}}
\providecommand{\nsplus}{\mathbin{+}}
\providecommand{\nssless}{<}
\providecommand{\nsless}{<}
\providecommand{\PL}{\mathrm{PL}}
\providecommand{\Nec}{\mathord{\Box}}
\providecommand{\RK}{\mathrm{RK}}
\providecommand{\ST}{\mathrm{ST}}
\providecommand{\SSubst}[3]{#1[#2/#3]}
\providecommand{\subst}[3]{#1[#2/#3]}
\providecommand{\tagfalse}{\mathrm{F}}
\providecommand{\cutr}{\mathrm{Cut}}
\providecommand{\cardneq}{\neq}
\providecommand{\issucc}[1]{\operatorname{succ}(#1)}
\providecommand{\tup}[1]{\langle #1 \rangle}
\providecommand{\Var}{\mathrm{Var}}
\providecommand{\Bew}{\operatorname{Bew}}
\providecommand{\Con}{\operatorname{Con}}
\providecommand{\Cons}{\operatorname{Cons}}
\providecommand{\pDeduce}{\mathrm{D}}
\providecommand{\psat}{\models}
\providecommand{\den}[1]{\llbracket #1 \rrbracket}
\providecommand{\Sem}[2]{\llbracket #1 \rrbracket_{#2}}
\providecommand{\DomainOf}[1]{\lvert #1 \rvert}
\providecommand{\Form}{\mathrm{Form}}
\providecommand{\Terms}{\mathrm{Term}}
\providecommand{\Atom}{\mathrm{Atom}}
\providecommand{\FVars}[1]{\operatorname{FV}(#1)}
\providecommand{\fv}[1]{\operatorname{fv}(#1)}
\providecommand{\tuplep}[1]{\langle #1 \rangle}
\providecommand{\restrict}[2]{#1 \upharpoonright #2}
\providecommand{\sto}{\rightarrow}
\providecommand{\FmlaSet}[1]{\mathrm{Fmla}_{#1}}
\providecommand{\Thm}{\mathrm{Thm}}
\providecommand{\Eval}{\operatorname{val}}
\providecommand{\blank}{\sqcup}
\providecommand{\blankTape}{\sqcup}
\providecommand{\Consis}{\operatorname{Con}}
\providecommand{\drs}{\mathrm{d}}
\providecommand{\VRule}{\mid}
\providecommand{\Bet}{\mathrm{B}}
\providecommand{\Seg}[1]{\mathrm{Seg}_{#1}}
\providecommand{\Enums}{\mathrm{En}}
\providecommand{\Z}{\mathbb{Z}}
\providecommand{\N}{\mathbb{N}}
\providecommand{\abs}[1]{\lvert #1 \rvert}
\providecommand{\floor}[1]{\lfloor #1 \rfloor}
\providecommand{\ceil}[1]{\lceil #1 \rceil}

% ===================== OLP compatibility layer (4) =====================
% --- fourth compatibility block: harvested from the per-candidate compile logs ---
\providecommand{\beredpar}{\mathrel{\rightarrow_\beta}}
\makeatletter
\def\bered{\@ifnextchar[{\bered@i}{\mathrel{\rightarrow_\beta}}}
\def\bered@i[#1]{\mathrel{\rightarrow_\beta}}
\def\bred{\@ifnextchar[{\bred@i}{\mathrel{\rightarrow_\beta}}}
\def\bred@i[#1]{\mathrel{\rightarrow_\beta}}
\makeatother
\providecommand{\beredone}{\mathrel{=_\beta}}
\providecommand{\bredone}{\mathrel{=_\beta}}
\providecommand{\xredone}{\mathrel{=_\beta}}
\providecommand{\aconvone}{\mathrel{=_\alpha}}
\providecommand{\cardless}{<}
\providecommand{\tagrefs}[1]{#1}
\providecommand{\ordplus}{\mathbin{+}}
\providecommand{\funrestrictionto}[2]{#1 \upharpoonright #2}
\providecommand{\fundefined}{\mathord{\uparrow}}
\providecommand{\Trm}{\mathrm{Trm}}
\providecommand{\OProv}{\operatorname{Prov}}
\providecommand{\OPrf}{\operatorname{Prf}}
\providecommand{\aeq}{\mathrel{=}}
\providecommand{\mClass}[1]{[#1]}
\providecommand{\equivclass}[1]{[#1]}
\providecommand{\pheight}{\operatorname{ht}}
\providecommand{\phy}{\operatorname{ht}}
\providecommand{\Cut}{\mathrm{Cut}}
\providecommand{\QuantRank}{\operatorname{qr}}
\providecommand{\elemequiv}{\mathrel{\equiv}}
\providecommand{\PIso}{\mathrm{PIso}}
\providecommand{\Theory}{\mathrm{Th}}
\providecommand{\supstrict}{\subsetneq}
\providecommand{\trcl}{\operatorname{trcl}}
\providecommand{\cref}[1]{\ref{#1}}
\providecommand{\Discharge}{\mathrm{Dis}}
\providecommand{\Weakening}{\mathrm{Wk}}
\providecommand{\ZFminus}{\mathrm{ZF}^{-}}
\providecommand{\RK}{\mathrm{RK}}
\makeatletter
\renewcommand{\Dual}{\@ifnextchar[{\Dual@i}{\ifmmode\mathrm{d}\else$\mathrm{d}$\fi}}
\def\Dual@i[#1]{\ifmmode\mathrm{d}[#1]\else$\mathrm{d}[#1]$\fi}
\makeatother
\providecommand{\RightR}{\mathrm{R}}
\providecommand{\LeftR}{\mathrm{L}}
\providecommand{\Intro}{\mathrm{I}}
\providecommand{\Axiom}{\mathrm{Ax}}
\providecommand{\formula}[1]{#1}
\providecommand{\stared}{\ast}
\providecommand{\oset}[2]{\{\, #1 \mid #2 \,\}}
\providecommand{\setof}[2]{\{\, #1 \mid #2 \,\}}
\providecommand{\pow}[1]{\mathcal{P}(#1)}
\providecommand{\ordpair}[2]{\langle #1, #2 \rangle}
\providecommand{\ordtriple}[3]{\langle #1, #2, #3 \rangle}
\providecommand{\union}{\cup}
\providecommand{\inter}{\cap}
\providecommand{\seq}[1]{\langle #1 \rangle}
\providecommand{\SuchThat}{\mid}
\providecommand{\where}{\mid}
\providecommand{\TVar}{\mathrm{TVar}}
\providecommand{\Bool}{\mathrm{Bool}}
\providecommand{\Val}{\operatorname{Val}}
\providecommand{\FV}{\operatorname{FV}}
\providecommand{\TL}[1]{\mathrm{TL}}
\providecommand{\xstared}{\ast}

% ===================== OLP compatibility layer (5) =====================
% --- fifth compatibility block: rescues the quarantined OLP candidates ---
\providecommand{\ident}{\operatorname{id}}
\providecommand{\setrank}{\operatorname{rank}}
\providecommand{\cardfont}{\mathfrak{c}}
\providecommand{\Complement}{\overline}
\providecommand{\QR}{\mathrm{QR}}
\providecommand{\becd}{\ensuremath{\ulcorner}}
\providecommand{\disjointsum}{\mathbin{\oplus}}
\providecommand{\cardeq}[2]{#1 = #2}
\providecommand{\xred}{\mathrel{\rightarrow}}
\providecommand{\cardless}{<}
\providecommand{\cardplus}{\mathbin{+}}
\providecommand{\cardtimes}{\mathbin{\cdot}}
\providecommand{\cardexpo}{\mathbin{\hat{\ }}}
\providecommand{\cardneq}{\neq}
\providecommand{\cardle}{\leq}
\providecommand{\cardnless}{\not<}
% --- sixth compatibility block: names reported by real compile logs ---
\providecommand{\ORefut}{\operatorname{Refut}}
\providecommand{\ORProv}{\operatorname{Prov}}
\providecommand{\cutrank}{\operatorname{cr}}
\providecommand{\Refut}{\operatorname{Refut}}
\providecommand{\DischargeRule}{\mathrm{Dis}}
\providecommand{\Contraction}{\mathrm{C}}
\providecommand{\bottomAlignProof}{\relax}
\providecommand{\FalseCl}{\bot}
\providecommand{\red}{\mathrel{\rightarrow}}
\providecommand{\varAssign}{\sigma}
\providecommand{\defiff}{\mathrel{\leftrightarrow}}
\providecommand{\FalseInt}{\bot}
\providecommand{\fn}{\operatorname{fn}}
\providecommand{\stared}{\ast}
\providecommand{\dotted}{\cdot}
\providecommand{\Gamma}{\Gamma}
% --- sixth compatibility block: names reported by real compile logs ---
\providecommand{\ORefut}{\operatorname{Refut}}
\providecommand{\ORProv}{\operatorname{Prov}}
\providecommand{\cutrank}{\operatorname{cr}}
\providecommand{\Refut}{\operatorname{Refut}}
\providecommand{\DischargeRule}{\mathrm{Dis}}
\providecommand{\Contraction}{\mathrm{C}}
\providecommand{\bottomAlignProof}{\relax}
\providecommand{\FalseCl}{\bot}
\providecommand{\red}{\mathrel{\rightarrow}}
\providecommand{\varAssign}{\sigma}
\providecommand{\defiff}{\mathrel{\leftrightarrow}}
\providecommand{\FalseInt}{\bot}
\providecommand{\fn}{\operatorname{fn}}
\providecommand{\stared}{\ast}
\providecommand{\dotted}{\cdot}
\providecommand{\Gamma}{\Gamma}
\begin{document}

% ==== OLP file: content/turing-machines/undecidability/decision-problem.tex ====
% Part: turing-machines
% Chapter: undecidability
% Section: unsolvability-decision-problem



\olfileid{tur}{und}{dec}
\olsection{The Decision Problem}

We say that first-order logic is \emph{decidable} iff there is an
effective method for determining whether or not a given !!{sentence}
is valid. As it turns out, there is no such method: the problem of
deciding validity of first-order sentences is unsolvable.

In order to establish this important negative result, we prove that
the decision problem cannot be solved by a Turing machine.  That is,
we show that there is no Turing machine which, whenever it is started
on a tape that contains a first-order !!{sentence}, eventually halts
and outputs either $1$ or~$0$ depending on whether the
!!{sentence} is valid or not. By the Church--Turing thesis, every
function which is computable is Turing computable. So if this
``validity function'' were effectively computable at all, it would be
Turing computable. If it isn't Turing computable, then, it also cannot
be effectively computable.

Our strategy for proving that the decision problem is unsolvable is to
reduce the halting problem to it.  This means the following: We have
proved that the function~$h(e,w)$ that halts with output~$1$ if the
Turing machine described by~$e$ halts on input~$w$ and outputs~$0$
otherwise, is not Turing computable.  We will show that if there were
a Turing machine that decides validity of first-order sentences, then
there is also Turing machine that computes~$h$.  Since $h$ cannot be
computed by a Turing machine, there cannot be a Turing machine that
decides validity either.

The first step in this strategy is to show that for every input~$w$
and a Turing machine~$M$, we can effectively describe !!a{sentence}
$!T(M, w)$ representing the instruction set of~$M$ and the input~$w$
and !!a{sentence}~$!E(M, w)$ expressing ``$M$ eventually halts'' such
that:
\begin{quote}
  $\Entails !T(M, w) \lif !E(M,w)$ iff $M$ halts for input~$w$.
\end{quote}
The bulk of our proof will consist in describing these sentences
$!T(M, w)$ and~$!E(M, w)$ and in verifying that $!T(M, w) \lif !E(M, w)$
is valid iff $M$~halts on input~$w$.


% ==== OLP file: content/turing-machines/undecidability/enumerating-tms.tex ====
% Part: turing-machines
% Chapter: undecidability
% Section: enumerating-tms



\olfileid{tur}{und}{enu}
\olsection{Enumerating Turing Machines}

\begin{explain}
We can show that the set of all Turing machines is !!{enumerable}.
This follows from the fact that each Turing machine can be finitely
described.  The set of states and the tape vocabulary are finite sets.
The transition function is a partial function from $Q \times \Sigma$
to $Q \times \Sigma \times \{\TMleft, \TMright, \TMstay\}$, and so
likewise can be specified by listing its values for the finitely many
argument pairs for which it is defined.

This is true as far as it goes, but there is a subtle difference. The
definition of Turing machines made no restriction on what !!{element}s
the set of states and tape alphabet can have. So, e.g., for every real
number, there technically is a Turing machine that uses that number as
a state. However, the \emph{behavior} of the Turing machine is
independent of which objects serve as states and vocabulary. Consider
the two Turing machines in \olref{fig:variants}.
\begin{figure}
\begin{center}
\begin{tikzpicture}[->,>=stealth',shorten >=1pt,auto,node distance=2.8cm,
                    semithick]
  \tikzstyle{every state}=[fill=none,draw=black,text=black]

  \node[initial,state]         (A)              {$q_0$};
  \node[state]         (B) [right of=A] {$q_1$};

  \path (A) edge [bend left] node {\TMtrans{\TMstroke}{\TMstroke}{\TMright}} (B)
        (B) edge [loop above] node {\TMtrans{\TMblank}{\TMblank}{\TMright}} (B)
            edge [bend left] node {\TMtrans{\TMstroke}{\TMstroke}{\TMright}} (A);
\end{tikzpicture}\\
\begin{tikzpicture}[->,>=stealth',shorten >=1pt,auto,node distance=2.8cm,
  semithick]
\tikzstyle{every state}=[fill=none,draw=black,text=black]

\node[initial,state]         (A)              {$s$};
\node[state]         (B) [right of=A] {$h$};

\path (A) edge [bend left] node {\TMtrans{A}{A}{\TMright}} (B)
(B) edge [loop above] node {\TMtrans{\TMblank}{\TMblank}{\TMright}} (B)
edge [bend left] node {\TMtrans{A}{A}{\TMright}} (A);
\end{tikzpicture}
\end{center}
\caption{Variants of the \emph{Even} machine}
\ollabel{fig:variants}
\end{figure}
These two diagrams correspond to two machines, $M$ with the tape
alphabet $\Sigma = \{\TMendtape,\TMblank,\TMstroke\}$ and set of
states $\{q_0,q_1\}$, and $M'$ with alphabet $\Sigma' =
\{\TMendtape,\TMblank,A\}$ and states $\{s,h\}$. But their
instructions are otherwise the same: $M$ will halt on a sequence of
$n$ $\TMstroke$'s iff $n$ is even, and $M'$ will halt on a sequence of
$n$ $A$'s iff $n$ is even. All we've done is rename $\TMstroke$
to~$A$, $q_0$ to~$s$, and $q_1$ to~$h$. This example generalizes: we
can think of Turing machines as the same as long as one results from
the other by such a renaming of symbols and states.  In fact, we can
simply think of the symbols and states of a Turing machine as positive
integers: instead of $\sigma_0$ think~$1$, instead of $\sigma_1$
think~$2$, etc.; $\TMendtape$ is~$1$, $\TMblank$ is~$2$, etc. In this
way, the \emph{Even} machine becomes the machine depicted in
\olref{fig:standard-even}.
\begin{figure}
\[\begin{tikzpicture}[->,>=stealth',shorten >=1pt,auto,node distance=2.8cm,
  semithick]
\tikzstyle{every state}=[fill=none,draw=black,text=black]

\node[initial,state]         (A)              {$1$};
\node[state]         (B) [right of=A] {$2$};

\path (A) edge [bend left] node {\TMtrans{3}{3}{\TMright}} (B)
(B) edge [loop above] node {\TMtrans{2}{2}{\TMright}} (B)
edge [bend left] node {\TMtrans{3}{3}{\TMright}} (A);
\end{tikzpicture}
\]
\caption{A standard \emph{Even} machine}
\ollabel{fig:standard-even}
\end{figure}
We might call a Turing machine with states and symbols that are
positive integers a \emph{standard} machine, and only consider
standard machines from now on.\footnote{The terminology ``standard
machine'' is not standard.}

We wanted to show that the set of Turing machines is !!{enumerable},
and with the above considerations in mind, it is enough to show that
the set of standard Turing machines is !!{enumerable}. Suppose we are
given a standard Turing machine $M = \tuple{Q, \Sigma, q_0, \delta}$.
How could we describe it using a finite string of positive integers?
We'll first list the number of states, the states themselves, the
number of symbols, the symbols themselves, and the starting state.
(Remember, all of these are positive integers, since $M$ is a standard
machine.)  What about~$\delta$? The set of possible arguments, i.e.,
pairs $\tuple{q,\sigma}$, is finite, since $Q$ and~$\Sigma$ are
finite. So the information in~$\delta$ is simply the finite list of
all $5$-tuples $\tuple{q, \sigma, q', \sigma', d}$ where
$\delta(q,\sigma) = \tuple{q', \sigma', D}$, and $d$ is a number that
codes the direction~$D$ (say, $1$ for~$\TMleft$, $2$ for~$\TMright$,
and $3$ for~$\TMstay$).

In this way, every standard Turing machine can be described by a
finite list of positive integers, i.e., as a sequence $s_M \in
(\PosInt)^*$. For instance, the standard \emph{Even} machine is coded
by the sequence
\[
2, \underbrace{1, 2}_Q, 3, \overbrace{1, 2, 3}^\Sigma, 1, \underbrace{1, 3, 2, 3, 2}_{\delta(1,3) = \tuple{2,3,R}}, 
\overbrace{2, 2, 2, 2, 2}^{\delta(2,2) = \tuple{2,2,R}},
\underbrace{2, 3, 1, 3, 2}_{\delta(2,3) = \tuple{1,3,R}}.
\]
\end{explain}

\begin{thm}
There are functions from $\Nat$ to~$\Nat$ which are not Turing
computable.
\end{thm}

\begin{proof}
We know that the set of finite sequences of positive
integers~$(\PosInt)^*$ is !!{enumerable}
(\cref{sfr:siz:zigzag:prob:posint-star}). This gives us that the set
of descriptions of standard Turing machines, as a subset
of~$(\PosInt)^*$, is itself enumerable.  Every Turing computable
function $\Nat$ to~$\Nat$ is computed by some (in fact, many) Turing
machines. By renaming its states and symbols to positive integers (in
particular, $\TMendtape$ as~$1$, $\TMblank$ as~$2$, and $\TMstroke$
as~$3$) we can see that every Turing computable function is computed
by a standard Turing machine. This means that the set of all Turing
computable functions from $\Nat$ to~$\Nat$ is also enumerable.

On the other hand, the set of all functions from $\Nat$ to~$\Nat$ is
not !!{enumerable} (\cref{sfr:siz:red:prob:nat-nat}). If all functions
were computable by some Turing machine, we could enumerate the set of
all functions by listing all the descriptions of Turing machines that
compute them. So there are some functions that are not Turing
computable. 
\end{proof}

\begin{prob}
  Can you think of a way to describe Turing machines that does not
  require that the states and alphabet symbols are explicitly listed?
  You may define your own notion of ``standard'' machine, but say
  something about why every Turing machine can be computed by a
  ``standard'' machine in your new sense.
\end{prob}


% ==== OLP file: content/turing-machines/undecidability/halting-problem.tex ====
% Part: computability
% Chapter: machines-computations
% Section: halting-problem



\olfileid{tur}{und}{hal}
\olsection{The Halting Problem}

\begin{explain}
Assume we have fixed some enumeration of Turing machine
descriptions.  Each Turing machine thus receives an \emph{index}: its
place in the enumeration $M_1$, $M_2$, $M_3$, \dots{} of Turing machine
descriptions.

We know that there must be non-Turing-computable functions: the set of
Turing machine descriptions---and hence the set of Turing machines---is
!!{enumerable}, but the set of all functions from $\Nat$ to $\Nat$ is
not. But we can find specific examples of non-computable functions as
well. One such function is the halting function.
\end{explain}

\begin{defn}[Halting function]
 The \emph{halting function}~$h$ is defined as
\[
h(e,n) =
\begin{cases}
  \text{0} & \text{if machine~$M_e$ does not halt for input $n$} \\
  \text{1} & \text{if machine~$M_e$ halts for input $n$}
\end{cases}
\]
\end{defn}

\begin{defn}[Halting problem]
The \emph{Halting Problem} is the problem of determining (for any $e$, $n$)
whether the Turing machine~$M_e$ halts for an input of~$n$ strokes.
\end{defn}

\begin{explain}
We show that~$h$ is not Turing-computable by showing that a related
function,~$s$, is not Turing-computable. This proof relies on the fact
that anything that can be computed by a Turing machine can be computed
by a disciplined Turing machine (\olref[mac][dis]{sec}), and the fact
that two Turing machines can be hooked together to create a single
machine (\olref[mac][cmb]{sec}).
\end{explain}

\begin{defn} The function~$s$ is defined as
\[
s(e) =
\begin{cases}
  \text{0} & \text{if machine~$M_e$ does not halt for input $e$} \\
  \text{1} & \text{if machine~$M_e$ halts for input $e$}
\end{cases}
\]
\end{defn}

\begin{lem}
The function~$s$ is not Turing computable.
\end{lem}

\begin{proof}
We suppose, for contradiction, that the function~$s$ is Turing
computable.  Then there would be a Turing machine~$S$ that
computes~$s$. We may assume, without loss of generality, that when $S$
halts, it does so while scanning the first square (i.e., that it is
disciplined).  This machine can be ``hooked up'' to another
machine~$J$, which halts if it is started on input~$0$ (i.e., if it
reads $\TMblank$ in the initial state while scanning the square to the
right of the end-of-tape symbol), and otherwise wanders off to the
right, never halting. $S \concat J$, the machine created by hooking
$S$ to~$J$, is a Turing machine, so it is $M_e$ for some~$e$ (i.e., it
appears somewhere in the enumeration). Start $M_e$ on an input of~$e$
$\TMstroke$s. There are two possibilities: either $M_e$ halts or it
does not halt.
\begin{enumerate}
\item Suppose $M_e$ halts for an input of $e$ $\TMstroke$s. Then $s(e)
  = 1$. So $S$, when started on~$e$, halts with a single $\TMstroke$
  as output on the tape.  Then $J$ starts with a $\TMstroke$ on the
  tape. In that case $J$ does not halt. But $M_e$ is the machine $S
  \concat J$, so it should do exactly what $S$ followed by $J$ would
  do (i.e., in this case, wander off to the right and never halt).  So
  $M_e$ cannot halt for an input of $e$ $\TMstroke$'s.

\item Now suppose $M_e$ does not halt for an input of $e$
  $\TMstroke$s.  Then $s(e) = 0$, and $S$, when started on input~$e$,
  halts with a blank tape.  $J$,~when started on a blank tape,
  immediately halts.  Again, $M_e$ does what $S$ followed by~$J$ would
  do, so $M_e$ must halt for an input of $e$ $\TMstroke$'s.
\end{enumerate}
In each case we arrive at a contradiction with our assumption. This
shows there cannot be a Turing machine~$S$: $s$~is not Turing
computable.
\end{proof}

\begin{thm}[Unsolvability of the Halting Problem]
\ollabel{thm:halting-problem} The halting problem is unsolvable, i.e.,
the function~$h$ is not Turing computable.
\end{thm}

\begin{proof}
Suppose $h$ were Turing computable, say, by a Turing machine~$H$. We
could use $H$ to build a Turing machine that computes~$s$: First, make
a copy of the input (separated by a~$\TMblank$ symbol). Then move back
to the beginning, and run~$H$.  We can clearly make a machine that
does the former (see \cref{tur:mac:dis:prob:copier}), and if $H$
existed, we would be able to ``hook it up'' to such a copier machine
to get a new machine which would determine if $M_e$ halts on
input~$e$, i.e., computes~$s$. But we've already shown that no such
machine can exist. Hence, $h$~is also not Turing computable.
\end{proof}

\begin{prob}
The Three Halting (3-Halt) problem is the problem of giving a decision
procedure to determine whether or not an arbitrarily chosen Turing
Machine halts for an input of three $\TMstroke$'s on an otherwise
blank tape. Prove that the 3-Halt problem is unsolvable.
\end{prob}

\begin{prob}
Show that if the halting problem is solvable for Turing machine and
input pairs $M_e$ and~$n$ where $e \neq n$, then it is also solvable
for the cases where $e = n$.
\end{prob}

\begin{prob}
We proved that the halting problem is unsolvable if the input is a
number~$e$, which identifies a Turing machine~$M_e$ via an enumeration
of all Turing machines.  What if we allow the description of Turing
machines from \olref[tur][und][enu]{sec} directly as input?  Can there
be a Turing machine which decides the halting problem but takes as
input descriptions of Turing machines rather than indices?  Explain
why or why not.
\end{prob}

\begin{prob} Show that the \emph{partial} function~$s'$ is defined as
  \[
  s'(e) =
  \begin{cases}
    \text{1} & \text{if machine~$M_e$ halts for input $e$}\\
    \text{undefined} & \text{if machine~$M_e$ does not halt for input $e$}
  \end{cases}
  \]
  \emph{is} Turing computable.
\end{prob}


% ==== OLP file: content/turing-machines/undecidability/introduction.tex ====
% Part: turing-machines 
% Chapter: undecidability
% Section: introduction



\olfileid{tur}{und}{int} 
\olsection{Introduction}

It might seem obvious that not every function, even every arithmetical
function, can be computable. There are just too many, whose behavior
is too complicated.  Functions defined from the decay of radioactive
particles, for instance, or other chaotic or random behavior. Suppose
we start counting 1-second intervals from a given time, and define the
function $f(n)$ as the number of particles in the universe that decay
in the $n$-th 1-second interval after that initial moment.  This seems
like a candidate for a function we cannot ever hope to compute.

But it is one thing to not be able to imagine how one would compute
such functions, and quite another to actually prove that they are
uncomputable.  In fact, even functions that seem hopelessly
complicated may, in an abstract sense, be computable.  For instance,
suppose the universe is finite in time---some day, in the very distant
future the universe will contract into a single point, as some
cosmological theories predict. Then there is only a finite (but
incredibly large) number of seconds from that initial moment for which
$f(n)$ is defined.  And any function which is defined for only finitely
many inputs is computable: we could list the outputs in one big table,
or code it in one very big Turing machine state transition diagram.

We are often interested in special cases of functions whose values give
the answers to yes/no questions.  For instance, the question ``is $n$
a prime number?'' is associated with the function
\[
\fn{isprime}(n) = \begin{cases}
  1 & \text{if $n$ is prime}\\
  0 & \text{otherwise.}
  \end{cases}
\]
We say that a yes/no question can be \emph{effectively decided}, if
the associated $1/0$-valued function is effectively computable.

To prove mathematically that there are functions which cannot be
effectively computed, or problems that cannot effectively decided, it
is essential to fix a specific model of computation, and show
that there are functions it cannot compute or problems it cannot
decide.  We can show, for instance, that not every function can be
computed by Turing machines, and not every problem can be decided by
Turing machines.  We can then appeal to the Church--Turing thesis to
conclude that not only are Turing machines not powerful enough to
compute every function, but no effective procedure can.

The key to proving such negative results is the fact that we can
assign numbers to Turing machines themselves.  The easiest way to do
this is to enumerate them, perhaps by fixing a specific way to write
down Turing machines and their programs, and then listing them in a
systematic fashion.  Once we see that this can be done, then the
existence of Turing-uncomputable functions follows by simple
cardinality considerations: the set of functions from $\Nat$ to~$\Nat$ (in
fact, even just from $\Nat$ to $\{0, 1\}$) are !!{nonenumerable}, but
since we can enumerate all the Turing machines, the set of Turing-computable
functions is only !!{denumerable}.

We can also define \emph{specific} functions and problems which we can
prove to be uncomputable and undecidable, respectively.  One such
problem is the so-called \emph{Halting Problem.} Turing machines can
be finitely described by listing their instructions.  Such a
description of a Turing machine, i.e., a Turing machine program, can
of course be used as input to another Turing machine.  So we can
consider Turing machines that decide questions about other Turing
machines.  One particularly interesting question is this: ``Does the
given Turing machine eventually halt when started on input~$n$?''  It
would be nice if there were a Turing machine that could decide this
question: think of it as a quality-control Turing machine which
ensures that Turing machines don't get caught in infinite loops and
such.  The interesting fact, which Turing proved, is that there cannot
be such a Turing machine. There cannot be a single Turing machine
which, when started on input consisting of a description of a Turing
machine $M$ and some number~$n$, will always halt with either output
$1$ or $0$ according to whether $M$ machine would have halted when
started on input $n$ or not.

Once we have examples of specific undecidable problems we can use them
to show that other problems are undecidable, too.  For instance, one
celebrated undecidable problem is the question, ``Is the first-order
!!{formula}~$!A$ valid?''.  There is no Turing machine which, given as
input a first-order !!{formula}~$!A$, is guaranteed to halt with
output $1$ or $0$ according to whether $!A$ is valid or not.
Historically, the question of finding a procedure to effectively solve
this problem was called simply ``the'' decision problem; and so we say
that the decision problem is unsolvable.  Turing and Church proved this
result independently at around the same time, so it is also called the
Church--Turing Theorem.


% ==== OLP file: content/turing-machines/undecidability/representing-tms.tex ====
% Part: turing-machines
% Chapter: undecidability
% Section: representing-tms



\olfileid{tur}{und}{rep}
\olsection{Representing Turing Machines}

\begin{explain}
In order to represent Turing machines and their behavior by !!a{sentence}
of first-order logic, we have to define a suitable language. The
language consists of two parts: !!{predicate}s for describing
configurations of the machine, and expressions for numbering execution
steps (``moments'') and positions on the tape.

We introduce two kinds of !!{predicate}s, both of them 2-place: For
each state~$q$, !!a{predicate}~$\Obj Q_q$, and for each tape
symbol~$\sigma$, !!a{predicate}~$\Obj S_\sigma$. The former allow us
to describe the state of~$M$ and the position of its tape head, the
latter allow us to describe the contents of the tape.  

In order to express the positions of the tape head and the number of
steps executed, we need a way to express numbers. This is done using
!!a{constant}~$\Obj 0$, and a $1$-place function~$\prime$, the
successor function. By convention it is written \emph{after} its
argument (and we leave out the parentheses).

For each number $n$ there is a canonical term~$\num{n}$, the
\emph{numeral} for~$n$, which represents it in~$\Lang L_M$. $\num{0}$
is $\Obj 0$, $\num{1}$ is $\Obj 0'$, $\num{2}$ is $\Obj 0''$, and so
on. More formally:
\begin{align*}
\num{0} & = \Obj 0 \\
\num{n+1} &= \num{n}'
\end{align*}
The term $\num{0}$, i.e., $\Obj 0$ names the leftmost position on the
tape as well as the time before the first execution step (the initial
configuration). The term $\num{1}$, i.e., $\Obj 0'$ names the square
to the right of the leftmost square, and the time after the first
execution step, and so on.

We also introduce !!a{predicate}~$<$ to express both the ordering of
tape positions (when it means ``to the left of'') and execution steps
(then it means ``before'').

Once we have the language in place, we list the ``axioms'' of $!T(M,
w)$, i.e., the !!{sentence}s which, taken together, describe the
behavior of~$M$ when run on input~$w$.  There will be !!{sentence}s
which lay down conditions on $\Obj 0$, $\prime$, and $<$,
!!{sentence}s that describes the input configuration, and
!!{sentence}s that describe what the configuration of $M$ is after it
executes a particular instruction.
\end{explain}

\begin{defn}
  \ollabel{defn:tm-descr}
Given a Turing machine $M = \tuple{Q, \Sigma, q_0, \delta}$, the
language~$\Lang L_M$ consists of:
\begin{enumerate}
\item A two-place !!{predicate} $\Obj Q_q(x, y)$ for every state~$q \in
  Q$.  Intuitively, $\Obj Q_q(\num{m}, \num{n})$ expresses ``after $n$
  steps, $M$ is in state~$q$ scanning the $m$th square.''
\item A two-place !!{predicate} $\Obj S_\sigma(x, y)$ for every
  symbol~$\sigma\in \Sigma$.  Intuitively, $\Obj S_\sigma(\num{m},
  \num{n})$ expresses ``after $n$ steps, the $m$th square contains
  symbol~$\sigma$.''
\item A !!{constant} $\Obj 0$
\item A one-place !!{function} $\prime$
\item A two-place !!{predicate} $<$
\end{enumerate}
\end{defn}

The !!{sentence}s describing the operation of the Turing machine~$M$ on
input $w = \sigma_{i_1}\dots\sigma_{i_k}$ are the following:
\begin{enumerate}
\item Axioms describing numbers and~$<$:
\begin{enumerate}
\item !!^a{sentence} that says that every number is less than its successor:
\[
\lforall[x][x < x']
\]
\item !!^a{sentence} that ensures that $<$ is transitive:
\[
\lforall[x][\lforall[y][\lforall[z][
      ((x < y \land y < z) \lif x < z)]]]
\]
\end{enumerate}
\item Axioms describing the input configuration:
\begin{enumerate}
\item After $0$~steps---before the machine starts---$M$ is in
  the initial state~$q_0$, scanning square~$1$:
\[
\Obj Q_{q_0}(\num{1}, \num{0})
\]
\item The first $k+1$ squares contain the symbols $\TMendtape$,
  $\sigma_{i_1}$, \dots, $\sigma_{i_k}$:
\[
\Obj S_\TMendtape(\num{0}, \num{0}) \land
\Obj S_{\sigma_{i_1}}(\num{1}, \num{0}) \land
\dots \land
\Obj S_{\sigma_{i_k}}(\num{k}, \num{0})
\]
\item Otherwise, the tape is empty:
\[
\lforall[x][(\num{k} < x \lif \Obj S_\TMblank(x, \num{0}))]
\]
\end{enumerate}
\item Axioms describing the transition from one configuration to
  the next:

For the following, let $!A(x, y)$ be the conjunction of all !!{sentence}s
of the form
\[
\lforall[z][
  (((z < x \lor x < z) \land \Obj S_\sigma(z, y))
  \lif \Obj S_\sigma(z, y'))]
\]
where $\sigma \in \Sigma$.  We use $!A(\num{m},\num{n})$ to express
``other than at square~$m$, the tape after $n+1$ steps is the same as
after $n$ steps.''
\begin{enumerate}
\item \ollabel{rep-right} For every instruction $\delta(q_i, \sigma) =
  \tuple{q_j, \sigma', \TMright}$, the !!{sentence}:
\begin{align*}
& \lforall[x][\lforall[y][(
   (\Obj Q_{q_i}(x, y) \land \Obj S_{\sigma}(x, y)) \lif {}]] \\
&\qquad   (\Obj Q_{q_j}(x', y') \land \Obj S_{\sigma'}(x, y') \land
!A(x, y)))
\end{align*}
This says that if, after~$y$ steps, the machine is in state~$q_i$
scanning square~$x$ which contains symbol~$\sigma$, then after $y+1$
steps it is scanning square~$x+1$, is in state~$q_j$, square~$x$ now
contains~$\sigma'$, and every square other than~$x$ contains the
same symbol as it did after~$y$ steps.

\item \ollabel{rep-left} For every instruction $\delta(q_i, \sigma) =
  \tuple{q_j, \sigma', \TMleft}$, the !!{sentence}:
\begin{align*}
& \lforall[x][\lforall[y][
    ((\Obj Q_{q_i}(x', y) \land \Obj S_{\sigma}(x', y)) \lif {}]]\\
& \qquad   (\Obj Q_{q_j}(x, y') \land \Obj S_{\sigma'}(x', y') \land
!A(x, y))) \land {}\\
& \lforall[y][((\Obj Q_{q_i}(\num{0}, y) \land \Obj S_{\sigma}(\num{0},
    y)) \lif {}]\\
& \qquad (\Obj Q_{q_j}(\num{0}, y') \land \Obj S_{\sigma'}(\num{0},
  y') \land !A(\num{0}, y)))
\end{align*}
Take a moment to think about how this works: now we don't start with
``if scanning square~$x$ \dots'' but: ``if scanning square $x+1$
\dots'' A move to the left means that in the next step the machine is
scanning square~$x$.  But the square that is written on is~$x+1$.  We
do it this way since we don't have subtraction or a predecessor
function.

Note that numbers of the form $x+1$ are $1$, $2$, \dots, i.e., this
doesn't cover the case where the machine is scanning square~$0$ and is
supposed to move left (which of course it can't---it just stays
put). That special case is covered by the second conjunction: it says
that if, after $y$ steps, the machine is scanning square~$0$ in state
$q_i$ and square~$0$ contains symbol~$\sigma$, then after $y+1$ steps
it's still scanning square~$0$, is now in state~$q_j$, the symbol on
square~$0$ is $\sigma'$, and the squares other than square~$0$ contain
the same symbols they contained ofter $y$~steps.
\item \ollabel{rep-stay} For every instruction $\delta(q_i, \sigma) =
  \tuple{q_j, \sigma', \TMstay}$, the !!{sentence}:
\begin{align*}
& \lforall[x][\lforall[y][(
   (\Obj Q_{q_i}(x, y) \land \Obj S_{\sigma}(x, y)) \lif {}]] \\
&\qquad   (\Obj Q_{q_j}(x, y') \land \Obj S_{\sigma'}(x, y') \land
!A(x, y)))
\end{align*}
\end{enumerate}
\end{enumerate}
Let $!T(M, w)$ be the conjunction of all the above !!{sentence}s for Turing
machine~$M$ and input~$w$.

In order to express that~$M$ eventually halts, we have to find a
!!{sentence} that says ``after some number of steps, the transition
function will be undefined.''  Let $X$~be the set of all pairs
$\tuple{q, \sigma}$ such that~$\delta(q, \sigma)$ is undefined.  Let
$!E(M, w)$ then be the !!{sentence}
\[
\lexists[x][\lexists[y][(\bigvee_{\tuple{q, \sigma} \in
      X}(\Obj Q_q(x, y) \land \Obj S_\sigma(x, y)))]]
\]

If we use a Turing machine with a designated halting state~$h$, it
is even easier: then the !!{sentence}~$!E(M, w)$
\[
\lexists[x][\lexists[y][\Obj Q_h(x, y)]]
\]
expresses that the machine eventually halts.

\begin{prop}
\ollabel{prop:mlessk}
If $m < k$, then $!T(M, w) \Entails \num{m} < \num{k}$
\end{prop}

\begin{proof}
Exercise.
\end{proof}

\begin{prob}
Prove \olref[tur][und][rep]{prop:mlessk}.
(Hint: use induction on $k-m$).
\end{prob}


% ==== OLP file: content/turing-machines/undecidability/trakhtenbrot.tex ====
% Part: turing-machines
% Chapter: undecidability
% Section: trakhtenbrot



\olfileid{tur}{und}{tra}
\olsection{Trakhtenbrot's Theorem}

\begin{explain}
  In \olref[rep]{sec} we defined !!{sentence}s $!T(M,w)$ and~$!E(M,w)$
  for a Turing machine~$M$ and input string~$w$. Then we showed in
  \olref[ver]{lem:valid-if-halt} and \olref[ver]{lem:halt-if-valid}
  that $!T(M,w) \lif !E(M,w)$ is valid iff $M$, started on input~$w$,
  eventually halts. Since the Halting Problem is undecidable, this
  implies that validity and satisfiability of !!{sentence}s of
  first-order logic is undecidable
  (\cref{tur:und:uns:thm:decision-prob,tur:und:uns:cor:undecidable-sat}).

  But validity and satisfiability of sentences is defined for
  arbitrary !!{structure}s, finite or infinite. You might suspect that
  it is easier to decide if !!a{sentence} is satisfiable in a finite
  !!{structure} (or valid in all finite !!{structure}s). We can adapt
  the proof of the unsolvability of the decision problem so that it
  shows this is not the case.

  First, if you go back to the proof of
  \olref[ver]{lem:halt-if-valid}, you'll see that what we did there is
  produce a model~$\Struct M$ of~$!T(M,w)$ which describes exactly
  what machine~$M$ does when started on input~$w$.  The domain of that
  model was~$\Nat$, i.e., infinite. But if $M$ actually halts on
  input~$w$, we can build a finite model~$\Struct M'$ in the same way.
  Suppose $M$ started on input~$w$ halts after~$k$ steps. Take as
  domain~$\Domain{M'}$ the set $\{0, \dots, n\}$, where $n$ is the
  larger of~$k$ and the length of~$w$, and let
  \[
    \Assign{\prime}{M'}(x) = 
    \begin{cases}
      x + 1 &\text{if $x < n$}\\
      n &\text{otherwise,}
    \end{cases}
  \]
  and $\tuple{x,y} \in \Assign{<}{M'}$ iff $x < y$ or $x = y = n$.
  Otherwise $\Struct{M'}$ is defined just like~$\Struct{M}$. By the
  definition of~$\Struct{M'}$, just like in the proof of
  \olref[ver]{lem:halt-if-valid}, $\Sat{M'}{!T(M,w)}$.  And since we
  assumed that $M$ halts on input~$w$, $\Sat{M'}{!E(M,w)}$. So,
  $\Struct{M'}$ is a finite model of~$!T(M,w) \land !E(M,w)$ (note
  that we've replaced $\lif$ with~$\land$).
  
  We are halfway to a proof: we've shown that if $M$ halts on
  input~$w$, then $!T(M,w) \land !E(M,w)$ has a finite model.
  Unfortunately, the converse of this does not hold, i.e., there are
  Turing machines that don't halt on some input~$w$, but $!T(M,w)
  \land !E(M,w)$ still has a finite model. For instance, consider the
  machine~$M$ with the single state $q_0$ and instruction
  $\delta(q_0,\TMblank) = \tuple{q_0,\TMblank,\TMstay}$. Started on
  empty input~$w = \emptyseq$, this machine never halts: it is in an
  infinite loop, but does not change the tape or move the head. All
  configurations are the same (same state, same head position, same
  tape contents). We can define a finite !!{structure}~$\Struct{M''}$
  that satisfies $!T(M,\emptyseq) \land !E(M,\emptyseq)$ (exercise).
  We can, however, change~$!T(M,w)$ in a suitable way so that such
  !!{structure}s are ruled out.

  \begin{prob}
    Let $M$ be a Turing machine with the single state $q_0$ and single
    instruction $\delta(q_0,\TMblank) = \tuple{q,\TMblank,\TMstay}$.
    Let $\Domain{M''} = \{0, 1, 2\}$, $\Assign{\prime}{M''}(0) =
    \Assign{\prime}{M'}(1) = 1$ and $\Assign{\prime}{M''}(2) = 2$, and
    $\Assign{<}{M''} = \{\tuple{0,1}, \tuple{1,1}, \tuple{2,2}\}$.
    Define $\Assign{\Obj Q_{q_0}}{M''}$, $\Assign{\Obj
    S_{\TMblank}}{M''}$, and $\Assign{\Obj S_{\TMendtape}}{M''}$ so
    that $!T(M,\emptyseq)$ and $!E(M,\emptyseq)$ become true and
    explain why they are. Hint: Observe that $\delta(q_0, \TMendtape)$
    is undefined. Ensure that
    \begin{align*}
    & \Obj Q_{q_0}(\num{1}, \num{n}) \land \Obj S_{\TMendtape}(\num{0}, \num{n})
      \land 
      \lforall[x][(\num{0} < x \lif \Obj S_\TMblank(x, \num{n}))] \text{\quad for all $n \in \Nat$}\\
    & \lexists[y][(\Obj Q_{q_0}(\num{0}, y) \land \Obj S_{\TMendtape}(\num{0}, y))]
    \end{align*}
    are both true in~$\Struct{M''}$.
  \end{prob}
\end{explain}

Consider the !!{sentence}s describing the operation of the Turing
machine~$M$ on input $w = \sigma_{i_1}\dots\sigma_{i_k}$:
\begin{enumerate}
\item Axioms describing numbers and~$<$ (just like in the definition
of~$!T(M,w)$ in \olref[rep]{sec}).
\item Axioms describing the input configuration: just like in the definition
of~$!T(M,w)$.
\item Axioms describing the transition from one configuration to
  the next:

For the following, let $!A(x, y)$ be as before, and let
\[
  !B(y) \ident \lforall[x][(x < y \lif \eq/[x][y])].
\]
\begin{enumerate}
\item \ollabel{rep-right} For every instruction $\delta(q_i, \sigma) =
  \tuple{q_j, \sigma', \TMright}$, the !!{sentence}:
\begin{align*}
& \lforall[x][\lforall[y][(
   (\Obj Q_{q_i}(x, y) \land \Obj S_{\sigma}(x, y)) \lif {}]] \\
&\qquad   (\Obj Q_{q_j}(x', y') \land \Obj S_{\sigma'}(x, y') \land
!A(x, y) \land !B(y')))
\end{align*}
\item \ollabel{rep-left} For every instruction $\delta(q_i, \sigma) =
  \tuple{q_j, \sigma', \TMleft}$, the !!{sentence}
\begin{align*}
& \lforall[x][\lforall[y][
    ((\Obj Q_{q_i}(x', y) \land \Obj S_{\sigma}(x', y)) \lif {}]]\\
& \qquad   (\Obj Q_{q_j}(x, y') \land \Obj S_{\sigma'}(x', y') \land
!A(x, y))) \land {}\\
& \lforall[y][((\Obj Q_{q_i}(\num{0}, y) \land \Obj S_{\sigma}(\num{0},
    y)) \lif {}]\\
& \qquad (\Obj Q_{q_j}(\num{0}, y') \land \Obj S_{\sigma'}(\num{0},
  y') \land !A(\num{0}, y) \land !B(y')))
\end{align*}
\item \ollabel{rep-stay} For every instruction $\delta(q_i, \sigma) =
  \tuple{q_j, \sigma', \TMstay}$, the !!{sentence}:
\begin{align*}
& \lforall[x][\lforall[y][(
   (\Obj Q_{q_i}(x, y) \land \Obj S_{\sigma}(x, y)) \lif {}]] \\
&\qquad (\Obj Q_{q_j}(x, y') \land \Obj S_{\sigma'}(x, y') \land
!A(x, y) \land !B(y')))
\end{align*}
\end{enumerate}
As you can see, the !!{sentence}s describing the transitions of~$M$ 
are the same as the corresponding !!{sentence}
in~$!T(M,w)$, except we add $!B(y')$ at the end. $!B(y')$~ensures
that the number $y'$ of the ``next'' configuration is different from
all previous numbers $\Obj 0$, $\Obj 0'$, \dots.
% \item Sentences that express consistency conditions:
% \begin{enumerate}
%   \item\ollabel{rep-con-symb}%
%   !!^a{sentence} that says that no square can have more than one
%   symbol on it at any given time:
%   \[\lforall[x][\lforall[y][(\Obj S_{\sigma}(x, y) \lif \lnot \Obj
%   S_{\sigma'}(x, y))]],\]
%   for every pair $\sigma \neq \sigma'$ of tape symbols of~$T$.
%   \item\ollabel{rep-con-state}%
%   !!^a{sentence} that says that $M$ cannot be in more than one state
%   at any given time:
%   \[\lforall[x][\lforall[y][(\Obj Q_{q}(x, y) \lif
%   \lforall[z][\lnot \Obj
%   Q_{q'}(z, y))]],\]
%   for every pair $q \neq q'$ of states of~$T$.
%   \item\ollabel{rep-con-tape}%
%   !!^a{sentence} that says that $M$ cannot be on more than one tape
%   square at any given time:
%   \[\lforall[x][\lforall[y][(\Obj Q_{q}(x, y) \lif
%   \lforall[z][\eq/[z][y]\lif \lnot \Obj
%   Q_{q'}(z, y))]],\]
%   for every pair $q$, $q'$ of states of~$T$.
% \end{enumerate}
\end{enumerate}
Let $!T'(M, w)$ be the conjunction of all the above !!{sentence}s for Turing
machine~$M$ and input~$w$.

\begin{lem}\ollabel{lem:halts-sat}
  If $M$ started on input~$w$ halts, then $!T'(M,w) \land !E(M,w)$ has
  a finite model.
\end{lem}

\begin{proof}
  Let $\Struct{M'}$ be as in the proof of
  \olref[ver]{lem:halt-if-valid}, except
  \begin{align*}
    \Domain{M'} & = \{0, \dots, n\},\\
    \Assign{\prime}{M'}(x) & = 
    \begin{cases}
      x + 1 &\text{if $x < n$}\\
      n &\text{otherwise,}
    \end{cases} \\
    \tuple{x,y} \in \Assign{<}{M'} &\text{iff $x < y$ or $x = y = n$,}
  \end{align*}
  where $n = \max(k,\len{w})$ and $k$~is the least number such that
  $M$ started on input~$w$ has halted after~$k$ steps. We leave the
  verification that $\Sat{M'}{!T'(M,w) \land E(M,w)}$ as an exercise.
\end{proof}

\begin{prob}
  Complete the proof of \olref[tur][und][tra]{lem:halts-sat} by
  proving that $\Sat{M'}{!T(M,w) \land E(M,w)}$.
\end{prob}

\begin{lem}\ollabel{lem:sat-halts}
  If $!T'(M,w) \land !E(M,w)$ has a finite model, then $M$ started on
  input~$w$ halts.
\end{lem}

\begin{proof}
  We show the contrapositive. Suppose that $M$ started on~$w$ does not
  halt. If $!T'(M,w) \land !E(M,w)$ has no model at all, we are done.
  So assume $\Struct{M}$ is a model of~$!T(M,w) \land !E(M, w)$. We
  have to show that it cannot be finite.
  
  We can prove, just like in \olref[ver]{lem:config}, that if~$M$,
  started on input~$w$, has not halted after~$n$ steps, then $!T'(M,w)
  \Entails !C(M, w, n) \land !B(\num{n})$. Since $M$ started on
  input~$w$ does not halt, $!T'(M,w) \Entails !C(M, w, n) \land
  !B(\num{n})$ for all~$n \in \Nat$. Note that by
  \olref[rep]{prop:mlessk}, $!T'(M,w) \Entails \num{k} < \num{n}$ for
  all~$k < n$. Also $!B(\num{n}) \Entails \num{k} < \num{n} \lif
  \eq/[\num{k}][\num{n}]$. So, $\Sat{M}{\eq/[\num{k}][\num{n}]}$ for
  all~$k < n$, i.e., the infinitely many terms~$\num{k}$ must all have
  different values in~$\Struct{M}$. But this requires that
  $\Domain{M}$ be infinite, so $\Struct{M}$ cannot be a finite model
  of~$!T'(M,w) \land !E(M, w)$.
\end{proof}

\begin{prob}
  Complete the proof of \olref[tur][und][tra]{lem:sat-halts} by
  proving that if~$M$, started on input~$w$, has not halted after~$n$
  steps, then $!T'(M,w) \Entails !B(\num{n})$.
\end{prob}

\begin{thm}[Trakhtenbrot's Theorem]
  \ollabel{thm:trakhtenbrodt}
  It is undecidable if an arbitrary !!{sentence} of first-order logic has
  a finite model (i.e., is finitely satisfiable).
\end{thm}

\begin{proof}
  Suppose there were a Turing machine~$F$ that decides the finite
  satisfiability problem. Then given any Turing machine~$M$ and
  input~$w$, we could compute the sentence~$!T'(M,w) \land !E(M,w)$,
  and use~$F$ to decide if it has a finite model. By
  \cref{tur:und:tra:lem:halts-sat,tur:und:tra:lem:sat-halts}, it does iff $M$
  started on input~$w$ halts. So we could use $F$ to solve the halting
  problem, which we know is unsolvable.
\end{proof}

\begin{cor}\ollabel{cor:fproof-incomp}%
  There can be no !!{derivation} system that is sound and complete for
  \emph{finite} validity, i.e., !!a{derivation} system which has $\Proves !B$
  iff $\Sat{M}{!B}$ for every finite !!{structure}~$\Struct{M}$.
\end{cor}

\begin{proof}
  Exercise.
\end{proof}

\begin{prob}
  Prove \olref[tur][und][tra]{cor:fproof-incomp}. Observe that $!B$ is
  satisfied in every finite !!{structure} iff $\lnot !B$ is not
  finitely satisfiable.  Explain why finite satisfiability is
  semi-decidable in the sense of
  \olref[tur][und][uns]{thm:valid-ce}. Use this to argue that if
  there were !!a{derivation} system for finite validity, then finite
  satisfiability would be decidable.
\end{prob}



% ==== OLP file: content/turing-machines/undecidability/undecidability.tex ====
% Part: turing-machines
% Chapter: undecidability



\olchapter{tur}{und}{Undecidability}

\olimport{introduction}

\olimport{enumerating-tms}

\olimport{universal-tm}

\olimport{halting-problem}

\olimport{decision-problem}

\olimport{representing-tms}

\olimport{verification}

\olimport{unsolvability-decision-problem}

\olimport{trakhtenbrot}

\OLEndChapterHook


% ==== OLP file: content/turing-machines/undecidability/universal-tm.tex ====
% Part: turing-machines
% Chapter: undecidability
% Section: universal-tm



\olfileid{tur}{und}{uni}
\olsection{Universal Turing Machines}

In \olref[enu]{sec} we discussed how every Turing machine can be
described by a finite sequence of integers. This sequence encodes the
states, alphabet, start state, and instructions of the Turing machine.
We also pointed out that the set of all of these descriptions is
!!{enumerable}. Since the set of such descriptions is !!{denumerable},
this means that there is !!a{surjective} function from~$\Nat$ to
these descriptions. Such !!a{surjective} function can be obtained, for
instance, using Cantor's zig-zag method.  It gives us a way of
enumerating all (descriptions) of Turing machines. If we fix one such
enumeration, it now makes sense to talk of the $1$st, $2$nd, \dots,
$e$th Turing machine.  These numbers are called \emph{indices}.

\begin{defn}
If $M$~is the $e$th Turing machine (in our fixed enumeration), we
say that $e$~is an \emph{index} of~$M$. We write $M_e$ for the $e$th
Turing machine.
\end{defn}

A machine may have more than one index, e.g., two descriptions of~$M$
may differ in the order in which we list its instructions, and these
different descriptions will have different indices.

Importantly, it is possible to give the enumeration of Turing machine
descriptions in such a way that we can effectively compute the
description of~$M$ from its index, and to effectively compute an index
of a machine~$M$ from its description.  By the Church--Turing thesis,
it is then possible to find a Turing machine which recovers the
description of the Turing machine with index~$e$ and writes the
corresponding description on its tape as output. The description would
be a sequence of blocks of~$\TMstroke$'s (representing the positive
integers in the sequence describing~$M_e$).

Given this, it now becomes natural to ask: what functions of Turing
machine indices are themselves computable by Turing machines? What
properties of Turing machine indices can be decided by Turing
machines?  An example: the function that maps an index~$e$ to the
number of states the Turing machine with index~$e$ has, is computable
by a Turing machine. Here's what such a Turing machine would do:
started on a tape containing a single block of $e$~$\TMstroke$'s, it
would first decode $e$ into its description. The description is now
represented by a sequence of blocks of~$\TMstroke$'s on the tape.
Since the first !!{element} in this sequence is the number of states.
So all that has to be done now is to erase everything but the first
block of $\TMstroke$'s and then halt.

A remarkable result is the following:

\begin{thm}\ollabel{thm:universal-tm} There is a \emph{universal
  Turing machine}~$U$ which, when started on input $\tuple{e,n}$ 
  \begin{enumerate}
    \item halts iff $M_e$ halts on input~$n$, and
    \item if $M_e$ halts with output $m$, so does~$U$.
  \end{enumerate}
  $U$ thus computes the function $f\colon \Nat \times \Nat \pto \Nat$
  given by $f(e,n) = m$ if $M_e$ started on input~$n$ halts with
  output~$m$, and undefined otherwise.
\end{thm}

\begin{proof}
  To actually produce~$U$ is basically impossible, since it is an
  extremely complicated machine. But we can describe in outline how it
  works, and then invoke the Church--Turing thesis.  When it starts,
  $U$'s tape contains a block of $e$ $\TMstroke$'s followed by a block
  of $n$~$\TMstroke$'s. It first ``decodes'' the index~$e$ to the
  right of the input~$n$. This produces a list of numbers (i.e.,
  blocks of $\TMstroke$'s separated by~$\TMblank$'s) that describes
  the instructions of machine~$M_e$. $U$ then writes the number of the
  start state of~$M_e$ and the number~$1$ on the tape to the right of
  the description of~$M_e$. (Again, these are represented in unary, as
  blocks of $\TMstroke$'s.) Next, it copies the input (block of
  $n$~$\TMstroke$'s) to the right---but it replaces each $\TMstroke$
  by a block of three $\TMstroke$'s (remember, the number of the $\TMstroke$
  symbol is~$3$, $1$ being the number of~$\TMendtape$ and $2$ being
  the number of~$\TMblank$). At the left end of this sequence of blocks
  (separated by $\TMblank$ symbols on the tape of~$U$), it writes a
  single~$\TMstroke$, the code for~$\TMendtape$.

  $U$ now has on its tape: the index~$e$, the number~$n$, the code
  number of the start state (the ``current state''), the number of the
  initial head position~$1$ (the ``current head position''), and the
  initial contents of the ``tape'' (a sequence of blocks
  of~$\TMstroke$'s representing the code numbers of the symbols
  of~$M_e$---the ``symbols''---separated by~$\TMblank$'s).

  It now simulates what $M_e$ would do if started on input~$n$, by
  doing the following:
  \begin{enumerate}
    \item Find the number~$k$ of the ``current head position'' (at the
    beginning, that's~$1$),
    \item Move to the $k$th block in the ``tape'' to see what the
    ``symbol'' there is,
    \item\ollabel{find-inst}%
    Find the instruction matching the current ``state'' and
    ``symbol,''
    \item Move back to the $k$th block on the ``tape'' and replace the
    ``symbol'' there with the code number of the symbol $M_e$ would
    write,
    \item Move the head to where it records the current ``state'' and
    replace the number there with the number of the new state,
    \item Move to the place where it records the ``tape position'' and
    erase a~$\TMstroke$ or add a~$\TMstroke$ (if the instruction says
    to move left or right, respectively).
    \item Repeat.\footnote{We're glossing over some subtle
    difficulties here. E.g., $U$~may need some extra space when it
    increases the counter where it keeps track of the ``current head
    position''---in that case it will have to move the entire ``tape''
    to the right.}
  \end{enumerate}
  If $M_e$ started on input~$n$ never halts, then $U$ also never
  halts, so its output is undefined.

  If in step~\olref{find-inst} it turns out that the description
  of~$M_e$ contains no instruction for the current
  ``state''/``symbol'' pair, then $M_e$ would halt. If this happens,
  $U$ erases the part of its tape to the left of the ``tape.'' For
  each block of three~$\TMstroke$'s (representing a~$\TMstroke$ on
  $M_e$'s tape), it writes a $\TMstroke$ on the left end of its own
  tape, and successively erases the ``tape.'' When this is done,
  $U$'s~tape contains a single block of~$\TMstroke$'s of length~$m$.
  
  If $U$ encounters something other than a block of
  three~$\TMstroke$'s on the ``tape,'' it immediately halts. Since
  $U$'s~tape in this case does not contain a single block
  of~$\TMstroke$'s, its output is not a natural number, i.e., $f(e,n)$
  is undefined in this case.
\end{proof}


% ==== OLP file: content/turing-machines/undecidability/unsolvability-decision-problem.tex ====
% Part: turing-machines 
% Chapter: undecidability 
% Section: unsolvability-decision-problem



\olfileid{tur}{und}{uns} 
\olsection{The Decision Problem is Unsolvable}

\begin{thm}
\ollabel{thm:decision-prob}
The decision problem is unsolvable: There is no Turing machine~$D$,
which when started on a tape that contains !!a{sentence}~$!B$ of
first-order logic as input, $D$~eventually halts, and outputs~$1$ iff
$!B$ is valid and $0$ otherwise.
\end{thm}

\begin{proof}
Suppose the decision problem were solvable, i.e., suppose there were a
Turing machine~$D$. Then we could solve the halting problem as
follows. We construct a Turing machine~$E$ that, given as input the
number~$e$ of Turing machine~$M_e$ and input~$w$, computes the
corresponding !!{sentence}~$!T(M_e, w) \lif !E(M_e, w)$ and halts,
scanning the leftmost square on the tape.  The machine $E \concat D$
would then, given input $e$ and $w$, first compute~$!T(M_e, w) \lif
!E(M_e, w)$ and then run the decision problem machine~$D$ on that
input.  $D$ halts with output~$1$ iff $!T(M_e, w) \lif !E(M_e, w)$ is
valid and outputs~$0$ otherwise. By \olref[ver]{lem:halt-if-valid} and
\olref[ver]{lem:valid-if-halt}, $!T(M_e, w) \lif !E(M_e, w)$ is valid
iff $M_e$ halts on input~$w$. Thus, $E\concat D$, given input $e$ and
$w$ halts with output~$1$ iff $M_e$ halts on input~$w$ and halts with
output~$0$ otherwise. In other words, $E \concat D$ would solve the
halting problem.  But we know, by \olref[hal]{thm:halting-problem},
that no such Turing machine can exist.
\end{proof}

\begin{cor}\ollabel{cor:undecidable-sat}%
It is undecidable if an arbitrary !!{sentence} of first-order logic is satisfiable.
\end{cor}

\begin{proof}
  Suppose satisfiability were decidable by a Turing machine~$S$. Then
  we could solve the decision problem as follows: Given
  !!a{sentence}~$B$ as input, move $!B$ to the right one square.
  Return to square~$1$ and write the symbol~$\lnot$.

  Now run the Turing machine~$S$. It eventually halts with output
  either $1$ (if $\lnot !B$ is satisfiable) or~$0$ (if $\lnot !B$ is
  unsatisfiable) on the tape. If there is a~$\TMstroke$ on square~$1$,
  erase it; if square~$1$ is empty, write a~$\TMstroke$, then halt.

  This Turing machine always halts, and its output is~$1$ iff $\lnot
  !B$ is unsatisfiable and $0$~otherwise. Since $!B$ is valid iff
  $\lnot !B$~is unsatisfiable, the machine outputs~$1$ iff $!B$ is
  valid, and $0$~otherwise, i.e., it would solve the decision problem.
\end{proof}

\begin{explain}
So there is no Turing machine which always gives a
correct ``yes'' or ``no'' answer to the question ``Is $!B$ a valid
!!{sentence} of first-order logic?'' However, there \emph{is} a Turing
machine that always gives a correct ``yes'' answer---but simply does
not halt if the answer is ``no.'' This follows from the soundness and
completeness theorem of first-order logic, and the fact that
!!{derivation}s can be effectively enumerated.
\end{explain}

\begin{thm}
  \ollabel{thm:valid-ce}%
  Validity of first-order !!{sentence}s is semi-decidable: There is a
  Turing machine~$E$, which when started on a tape that contains
  !!a{sentence}~$!B$ of first-order logic as input, $E$~eventually
  halts and outputs~$1$ iff $!B$ is valid, but does not halt
  otherwise.
\end{thm}

\begin{proof}
  All possible !!{derivation}s of first-order logic can be generated,
  one after another, by an effective algorithm.  The machine~$E$ does
  this, and when it finds !!a{derivation} that shows that $\Proves
  !B$, it halts with output~$1$. By the soundness theorem, if $E$
  halts with output~$1$, it's because~$\Entails !B$. By the
  completeness theorem, if $\Entails !B$ there is !!a{derivation} that
  shows that~$\Proves !B$. Since $E$ systematically generates all
  possible !!{derivation}s, it will eventually find one that
  shows~$\Proves !B$, so will eventually halt with output~$1$.
\end{proof}


% ==== OLP file: content/turing-machines/undecidability/verification.tex ====
% Part: turing-machines 
% Chapter: undecidability 
% Section: verification



\olfileid{tur}{und}{ver} 
\olsection{Verifying the Representation}

\begin{explain}
In order to verify that our representation works, we have to prove two
things. First, we have to show that if $M$ halts on input~$w$, then
$!T(M, w) \lif !E(M, w)$ is valid. Then, we have to show the converse,
i.e., that if $!T(M, w) \lif !E(M, w)$ is valid, then $M$ does in fact
eventually halt when run on input~$w$.

The strategy for proving these is very different. For the first
result, we have to show that !!a{sentence} of first-order logic
(namely, $!T(M, w) \lif !E(M, w)$) is valid. The easiest way to do
this is to give !!a{derivation}. Our proof is supposed to work for all
$M$ and $w$, though, so there isn't really a single !!{sentence} for
which we have to give !!a{derivation}, but infinitely many.  So the best
we can do is to prove by induction that, whatever $M$ and~$w$ look
like, and however many steps it takes $M$ to halt on input~$w$, there
will be !!a{derivation} of $!T(M, w) \lif !E(M, w)$.

Naturally, our induction will proceed on the number of steps $M$ takes
before it reaches a halting configuration. In our inductive proof,
we'll establish that for each step~$n$ of the run of $M$ on input~$w$,
$!T(M, w) \Entails !C(M, w, n)$, where $!C(M, w, n)$ correctly
describes the configuration of $M$ run on~$w$ after~$n$ steps. Now if
$M$ halts on input~$w$ after, say, $n$ steps, $!C(M, w, n)$ will
describe a halting configuration. We'll also show that $!C(M, w, n)
\Entails !E(M, w)$, whenever $!C(M, w, n)$ describes a halting
configuration.  So, if $M$ halts on input~$w$, then for some~$n$, $M$
will be in a halting configuration after $n$~steps. Hence, $!T(M, w)
\Entails !C(M, w, n)$ where $!C(M, w, n)$ describes a halting
configuration, and since in that case $!C(M, w, n) \Entails !E(M, w)$,
we get that $T(M, w) \Entails !E(M, w)$, i.e., that $\Entails !T(M, w)
\lif !E(M, w)$.

The strategy for the converse is very different. Here we assume that
$\Entails !T(M, w) \lif !E(M, w)$ and have to prove that $M$ halts on
input~$w$.  From the hypothesis we get that $!T(M, w) \Entails !E(M,
w)$, i.e., $!E(M, w)$ is true in every !!{structure} in which $!T(M,
w)$ is true. So we'll describe !!a{structure}~$\Struct{M}$ in which
$!T(M, w)$ is true: its domain will be $\Nat$, and the interpretation
of all the $\Obj Q_q$ and $\Obj S_\sigma$ will be given by the
configurations of~$M$ during a run on input~$w$.  So, e.g.,
$\Sat{M}{\Obj Q_q(\num{m}, \num{n})}$ iff $T$, when run on input~$w$
for $n$ steps, is in state~$q$ and scanning square~$m$.  Now since
$!T(M, w) \Entails !E(M, w)$ by hypothesis, and since $\Sat{M}{!T(M,
  w)}$ by construction, $\Sat{M}{!E(M, w)}$.  But $\Sat{M}{!E(M, w)}$
iff there is some $n \in \Domain{M} = \Nat$ so that $M$, run on
input~$w$, is in a halting configuration after~$n$ steps.
\end{explain}

\begin{defn} 
Let $!C(M, w, n)$ be the !!{sentence}
\[ 
\Obj Q_q(\num{m}, \num{n}) \land \Obj S_{\sigma_0}(\num{0}, \num{n})
\land \dots \land \Obj S_{\sigma_k}(\num{k}, \num{n}) \land
\lforall[x][(\num{k} < x \lif \Obj S_\TMblank(x, \num{n}))]
\] 
where $q$ is the state of $M$ at time~$n$, $M$ is scanning square~$m$
at time~$n$, square~$i$ contains symbol~$\sigma_i$ at time~$n$ for $0
\le i \le k$ and $k$ is the right-most non-blank square of the tape at
time~$0$, or the right-most square the tape head has visited after $n$
steps, whichever is greater.
\end{defn}

\begin{lem}\ollabel{lem:halt-config-implies-halt}
If $M$ run on input~$w$ is in a halting configuration after $n$ steps,
then $!C(M, w, n) \Entails !E(M, w)$.
\end{lem}

\begin{proof}
Suppose that $M$ halts for input~$w$ after $n$ steps.
There is some state~$q$, square~$m$, and symbol~$\sigma$ such that:
\begin{enumerate} 
\item After $n$ steps, $M$ is in state~$q$ scanning square~$m$ on
  which~$\sigma$ appears.
\item The transition function $\delta(q, \sigma)$ is undefined.
\end{enumerate}
$!C(M, w, n)$ is the description of this configuration and will
include the clauses $\Obj Q_{q}(\num{m}, \num{n})$ and $\Obj
S_{\sigma}(\num{m}, \num{n})$. These clauses together imply $!E(M,
w)$:
\[
\lexists[x][\lexists[y][(\bigvee_{\tuple{q, \sigma} \in
      X}(\Obj Q_q(x, y) \land \Obj S_\sigma(x, y)))]]
\]
since $\Obj Q_{q'}(\num{m}, \num{n}) \land S_{\sigma'}(\num{m},
\num{n}) \Entails \bigvee_{\tuple{q, \sigma} \in X} (\Obj Q_q(\num{m},
\num{n}) \land \Obj S_{\sigma}(\num{m}, \num{n}))$, as
$\tuple{q',\sigma'} \in X$. 
\end{proof}

\begin{explain}
So if $M$ halts for input $w$, then there is some~$n$ such that $!C(M,
w, n) \Entails !E(M,w)$.  We will now show that for any time~$n$,
$!T(M, w) \Entails !C(M, w, n)$.
\end{explain}

\begin{lem}
\ollabel{lem:config}
For each $n$, if $M$ has not halted after $n$ steps, $!T(M, w)
\Entails !C(M, w, n)$.
\end{lem}

\begin{proof}
Induction basis: If $n = 0$, then the conjuncts of $!C(M, w, 0)$ are
also conjuncts of $!T(M, w)$, so entailed by it.

Inductive hypothesis: If $M$ has not halted before the $n$th step,
then $!T(M,w) \Entails !C(M, w, n)$. We have to show that (unless
$!C(M, w, n)$ describes a halting configuration), $!T(M, w) \Entails
!C(M, w, n+1)$.

Suppose $n > 0$ and after $n$ steps, $M$ started on $w$ is in
state~$q$ scanning square~$m$. Since $M$ does not halt after~$n$
steps, there must be an instruction of one of the following three
forms in the program of~$M$:

\begin{enumerate} 
\item \ollabel{right} $\delta(q, \sigma) = \tuple{q', \sigma', \TMright}$

\item \ollabel{left} $\delta(q, \sigma) = \tuple{q', \sigma', \TMleft}$

\item \ollabel{stay} $\delta(q, \sigma) = \tuple{q', \sigma', \TMstay}$
\end{enumerate}

We will consider each of these three cases in turn. 

\begin{enumerate} 
\item Suppose there is an instruction of the form~\olref{right}.
  By \olref[rep]{defn:tm-descr}\olref[rep]{rep-right}, this means that
\begin{align*} 
& \lforall[x][\lforall[y][((\Obj Q_{q}(x,
  y) \land \Obj S_\sigma(x, y)) \lif {}]]\\
& \qquad (\Obj Q_{q'}(x',
  y') \land \Obj S_{\sigma'}(x, y') \land !A(x, y)))  
\intertext{is a conjunct of $!T(M,w)$. This entails the following
  !!{sentence} (universal instantiation, $\num{m}$ for~$x$ and
  $\num{n}$ for~$y$):}
& (\Obj Q_{q}(\num{m}, \num{n}) \land \Obj S_{\sigma}(\num{m},
\num{n})) \lif {}\\
& \qquad (\Obj Q_{q'}(\num{m}', \num{n}') \land
\Obj S_{\sigma'}(\num{m}, \num{n}') \land !A(\num{m}, \num{n})).
\intertext{By induction hypothesis, $!T(M, w) \Entails !C(M, w, n)$,
  i.e.,}
& \Obj Q_q(\num{m}, \num{n}) \land \Obj S_{\sigma_0}(\num{0}, \num{n})
\land \dots \land \Obj S_{\sigma_k}(\num{k}, \num{n}) \land \\
& \qquad\lforall[x][(\num{k} < x \lif \Obj S_\TMblank(x, \num{n}))]\\
\intertext{Since after $n$ steps, tape square~$m$ contains~$\sigma$,
  the corresponding conjunct is~$\Obj S_\sigma(\num{m}, \num{n})$,
  so this entails:}
& \Obj Q_{q}(\num{m}, \num{n}) \land \Obj S_{\sigma}(\num{m},
   \num{n})
\intertext{We now get}
& \Obj Q_{q'}(\num{m}', \num{n}') \land \Obj S_{\sigma'}(\num{m},
  \num{n}') \land {}\\
& \qquad\Obj S_{\sigma_0}(\num{0}, \num{n}') \land \dots \land
  \Obj S_{\sigma_k}(\num{k}, \num{n}') \land {}\\
& \qquad \lforall[x][(\num{k} < x \lif \Obj S_\TMblank(x, \num{n}'))]
\end{align*}
as follows: The first line comes directly from the consequent of the
preceding conditional, by modus ponens. Each conjunct in the middle
line---which excludes $S_{\sigma_m}(\num{m},\num{n}')$---follows from
the corresponding conjunct in~$!C(M, w, n)$ together with $!A(\num{m},
\num{n})$.

If $m < k$, $!T(M,w) \Proves \num{m} < \num {k}$
(\olref[rep]{prop:mlessk}) and by transitivity of~$<$, we have
$\lforall[x][(\num{k} < x \lif \num{m} < x)]$.  If $m = k$, then
$\lforall[x][(\num{k} < x \lif \num{m} < x)]$ by logic alone.  The
last line then follows from the corresponding conjunct in $!C(M, w,
n)$, $\lforall[x][(\num{k} < x \lif \num{m} < x)]$, and $!A(\num{m},
\num{n})$.  If $m<k$, this already is $!C(M, w, n+1)$.

Now suppose $m=k$. In that case, after $n+1$ steps, the tape head has
also visited square~$k+1$, which now is the right-most square
visited.  So $!C(M, w, n+1)$ has a new conjunct, $\Obj
S_\TMblank(\num{k}',\num{n}')$, and the last conjunct is
$\lforall[x][(\num{k}' < x \lif \Obj S_\TMblank(x, \num{n}'))]$. We
have to verify that these two !!{sentence}s are also implied.

We already have $\lforall[x][(\num{k} < x \lif \Obj S_\TMblank(x,
  \num{n}'))]$. In particular, this gives us $\num{k} < \num{k}' \lif
\Obj S_\TMblank(\num{k}', \num{n}')$. From the axiom $\lforall[x][x <
  x']$ we get $\num{k} < \num{k}'$. By modus ponens, $\Obj
S_\TMblank(\num{k}',\num{n}')$ follows.

Also, since $!T(M,w) \Proves \num{k} < \num{k}'$, the axiom for
transitivity of~$<$ gives us $\lforall[x][(\num{k}' < x \lif \Obj
  S_\TMblank(x, \num{n}'))]$. (We leave the verification of this as an
exercise.)

\item Suppose there is an instruction of the form~\olref{left}.
  Then, by \olref[rep]{defn:tm-descr}\olref[rep]{rep-left},
\begin{align*} 
& \lforall[x][\lforall[y][((\Obj Q_{q}(x', y) \land \Obj
    S_{\sigma}(x', y)) \lif {}]]\\
& \qquad (\Obj Q_{q'}(x, y') \land \Obj
  S_{\sigma'}(x', y') \land !A(x, y))) \land {}\\
& \lforall[y][((\Obj Q_{q_i}(\num{0}, y) \land \Obj S_{\sigma}(\num{0},
    y)) \lif {}]\\
& \qquad (\Obj Q_{q_j}(\num{0}, y') \land \Obj S_{\sigma'}(\num{0}, y')
  \land !A(\num{0}, y)))
\intertext{is a conjunct of $!T(M,w)$. If $m>0$, then let $l = m - 1$
  (i.e., $m = l+1$). The first conjunct of the above !!{sentence}
  entails the following:}
& (\Obj Q_{q}(\num{l}', \num{n}) \land \Obj S_{\sigma}(\num{l}', \num{n}))
\lif {} \\
& \qquad (\Obj Q_{q'}(\num{l}, \num{n}') \land \Obj S_{\sigma'}(\num{l}',
\num{n}') \land !A(\num{l}, \num{n}))
\intertext{Otherwise, let $l = m = 0$ and consider the following !!{sentence}
  entailed by the second conjunct:}
& ((\Obj Q_{q_i}(\num{0}, \num{n}) \land \Obj S_{\sigma}(\num{0}, \num{n})) \lif {}\\
& \qquad   (\Obj Q_{q_j}(\num{0}, \num{n}') \land \Obj S_{\sigma'}(\num{0}, \num{n}') \land
!A(\num{0}, \num{n}))) 
\intertext{Either sentence implies}
&  \Obj Q_{q'}(\num{l}, \num{n}') \land \Obj S_{\sigma'}(\num{m},
  \num{n}') \land {}\\
&\qquad  \Obj S_{\sigma_0}(\num{0}, \num{n}')\land \dots \land
  \Obj S_{\sigma_k}(\num{k}, \num{n}')  \land {} \\
& \qquad  \lforall[x][(\num{k} < x
  \lif \Obj S_\TMblank(x, \num{n}'))]
\end{align*}
as before. (Note that in the first case, $\num{l}' \ident \num{l+1}
\ident \num{m}$ and in the second case $\num{l} \ident \num{0}$.) But
this just is $!C(M, w, n+1)$.

\item Case \olref{stay} is left as an exercise.
\end{enumerate}
We have shown that for any~$n$, $!T(M, w) \Entails !C(M, w, n)$.
\end{proof}

\begin{prob}
Complete case~\olref[tur][und][ver]{stay} of the proof of
\olref[tur][und][ver]{lem:config}.
\end{prob}

\begin{prob}
Give !!a{derivation} of $\Obj S_{\sigma_i}(\num{i}, \num{n}')$ from
$\Obj S_{\sigma_i}(\num{i}, \num{n})$ and $!A(m, n)$ (assuming $i \neq
m$, i.e., either $i < m$ or $m < i$).
\end{prob}

\begin{prob}
Give !!a{derivation} of $\lforall[x][(\num{k}' < x \lif \Obj
  S_\TMblank(x, \num{n}'))]$ from $\lforall[x][(\num{k} < x \lif \Obj
  S_\TMblank(x, \num{n}'))]$, $\lforall[x][x < x']$, and
$\lforall[x][\lforall[y][\lforall[z][ ((x < y \land y < z) \lif x <
      z)]]]$.)
\end{prob}


\begin{lem}
\ollabel{lem:valid-if-halt}
If $M$ halts on input~$w$, then $!T(M, w) \lif
!E(M, w)$ is valid.
\end{lem}

\begin{proof}
By \olref{lem:config}, we know that, for any time~$n$, the
description~$!C(M, w, n)$ of the configuration of $M$ at time~$n$ is
entailed by~$!T(M, w)$.  Suppose $M$ halts after $k$ steps. At that
point, it will be scanning square~$m$, for some~$m \in \Nat$. Then
$!C(M, w, k)$ describes a halting configuration of~$M$, i.e., it
contains as conjuncts both $\Obj Q_q(\num{m}, \num{k})$ and $\Obj
S_\sigma(\num{m}, \num{k})$ with $\delta(q,\sigma)$ undefined.  Thus,
by \olref{lem:halt-config-implies-halt}, $!C(M, w, k) \Entails !E(M,
w)$. But since $!T(M, w) \Entails !C(M, w, k)$, we have $!T(M, w)
\Entails !E(M, w)$ and therefore $!T(M, w) \lif !E(M, w)$ is valid.
\end{proof}


\begin{explain} 
To complete the verification of our claim, we also have to
establish the reverse direction: if $!T(M, w) \lif !E(M, w)$ is valid, then
$M$ does in fact halt when started on input~$w$. 
\end{explain}

\begin{lem}
\ollabel{lem:halt-if-valid}
If $\Entails !T(M, w) \lif !E(M, w)$, then $M$ halts on input~$w$.
\end{lem}

\begin{proof}
Consider the $\Lang L_M$-!!{structure}~$\Struct M$ with
domain~$\Nat$ which interprets $\Obj 0$ as~$0$, $\prime$~as the successor
function, and $<$~as the less-than relation, and the predicates $\Obj Q_q$
and~$\Obj S_\sigma$ as follows:
\begin{align*}
  \Assign{\Obj Q_q}{M} & =
\Setabs{\tuple{m, n}}{\begin{array}{ll}\text{started on~$w$, after~$n$ steps,}\\ \text{$M$ is in state $q$
  scanning square~$m$}\end{array}} \\
\Assign{\Obj S_\sigma}{M} & = \Setabs{\tuple{m, n}}{\begin{array}{ll}
\text{started on~$w$, after $n$ steps,}\\ \text{square~$m$ of $M$ contains
  symbol~$\sigma$}\end{array}}
\end{align*}
In other words, we construct the !!{structure}~$\Struct{M}$ so that it
describes what $M$ started on input~$w$ actually does, step by step.
Clearly, $\Sat{M}{!T(M, w)}$. If $\Entails !T(M, w) \lif !E(M, w)$,
then also $\Sat{M}{!E(M, w)}$, i.e.,
\[
\Sat{M}{\lexists[x][\lexists[y][(\bigvee_{\tuple{q, \sigma} \in
      X}(\Obj Q_q(x, y) \land \Obj S_\sigma(x, y)))]]}.
\]
As $\Domain{M} = \Nat$, there must be $m$, $n \in \Nat$ so that
$\Sat{M}{\Obj Q_q(\num{m}, \num{n}) \land \Obj S_\sigma(\num{m},
\num{n})}$ for some~$q$ and~$\sigma$ such that $\delta(q, \sigma)$ is
undefined. By the definition of~$\Struct M$, this means that $M$
started on input~$w$ after~$n$ steps is in state~$q$ and reading
symbol~$\sigma$, and the transition function is undefined, i.e.,
$M$~has halted.
\end{proof}


